\chapter*{Conferencia IV}
\addcontentsline{toc}{chapter}{LECTURA I - El timbre y el carácter del violín}
SEÑORES: En este momento les presento una de las cuestiones más fascinantes relacionadas con el violín: la madera de la caja de resonancia. Para algunos estudiantes de violín, la cuestión del barniz resulta particularmente atractiva. Para otros, la del color puede serlo. El estudio del color, en sí mismo, constituye una rama del arte. Nadie debería creerse capaz de dominar el trabajo con el color del violín en poco tiempo. Tampoco se puede comprar el dominio del trabajo con el color por una cantidad determinada.
de efectivo.
[Tuve un amigo aficionado a la construcción de violines que disfrutaba de todas las ventajas de la riqueza y que perdió la vida en un segundo viaje por Europa en busca de colores para el violín. ¡Pobre hombre! Imaginó ingenuamente que el dinero podía comprar un talento que es un don o el resultado de una larga práctica. El sentido del color no es algo que se adquiera comercialmente. No tuvo la paciencia necesaria para practicar el trabajo con el color durante los años que requiere su desarrollo.]
Pero, a pesar del gran interés que despiertan el barniz y los colores, la madera de la caja de resonancia sigue siendo de suma importancia. Parece ser unánime la opinión de que la calidad de la madera de la caja de resonancia influye considerablemente en el timbre del violín.
Todo lo que contiene o rodea a la madera de caja de resonancia ha sido sometido a un escrutinio minucioso. En dicho escrutinio, la ciencia ha brindado escasa ayuda. Solo a la experimentación le debemos un conocimiento limitado de la misma. Por cierto, la ciencia ofrece una
Un dato de valor corroborativo en el registro de la conductividad de las ondas sonoras en el aire, los gases, el agua, los metales y las distintas variedades de madera. La botánica aporta un dato al constatar que los capilares del duramen se reducen a medida que el árbol envejece, por lo que las fibras del duramen poseen una mayor elasticidad.
[Para ser precisos, esta afirmación no incluye las fibras que se encuentran inmediatamente en el corazón, sino las que están cerca del corazón.]
Como dato de interés para el estudiante de violín, cito algunos puntos de la tabla de conductividad de ese admirable tratado sobre la teoría del sonido en relación con la música, de Pietro Blaserna, de la Real Universidad de Roma.
[Esta obra es completa en detalles y de dicción impecable, excepto en los casos en que el traductor, a lo largo del libro, con solo dos excepciones aisladas, usa erróneamente la palabra "nota" cuando significa
"tono"] Las cifras de Blaserna se dan en el sistema métrico, a saber: 

\begin{quote}
Velocidad del sonido en varios cuerpos:

Aire (32°F): 330 m/s

Cobre (68°F): 3556 m/s

Madera de acacia:
\begin{itemize}
\item A lo largo de las fibras: 4714 m/s
\item A través de los anillos: 1458 m/s
\item Con los anillos: 2954 m/s
\end{itemize}
\end{quote}

Estas cifras presentan algunos datos interesantes para el estudiante de violín. Así: (1) Las ondas sonoras a lo largo de la fibra del pino viajan a más de diez veces la velocidad que viajan en el aire. (2) En el cobre (y en todos los demás metales), a medida que aumenta la temperatura, la distancia que recorren las ondas sonoras disminuye.
[La conductividad de los metales, que disminuye a temperaturas elevadas, se introduce aquí para destacar que tanto las altas como las bajas temperaturas modifican considerablemente la conductividad de todos los cuerpos. Así, con el aumento de la temperatura del aire propio del verano, todos los sonidos, musicales o de otro tipo, solo pueden propagarse a distancias relativamente menores. El mismo fenómeno ocurre con las bajas temperaturas del invierno. Por lo tanto, la intensidad o la capacidad de transmisión de un violín no debe determinarse mediante una prueba realizada en temperaturas extremas.]
Pero una variedad de madera supera al pino en su capacidad de conducir el sonido. Esa madera, la acacia o canela, no está disponible para su uso como caja de resonancia. Por lo tanto, el pino se coloca al principio de la lista. Pero como existen varias especies del género Pinus, debemos elegir entre ellas. Es lamentable que no se indique la especie particular de pino empleada para realizar la grabación. Podemos suponer legítimamente que la muestra utilizada para la grabación se tomó cerca del corazón de un árbol maduro, porque, como se afirma en botánica, los tubos capilares del duramen en el árbol maduro disminuyen de calibre con la edad hasta que dejan de transportar savia. Por lo tanto, el duramen posee mayor densidad.
que la albura. La madera del duramen también posee una mayor elasticidad que la albura, y también mayor peso. Pero, una vez curada, la albura puede ser más difícil de cortar que el duramen. La mera dureza al cortar, hasta donde alcanza mi observación, es una cualidad perjudicial para la máxima calidad tonal. Al examinar el tocón nuevo de un árbol maduro, los crecimientos anuales recientes bajo la corteza aún no están bien definidos en fibra dura y tejido conectivo, sino que parecen más bien una masa homogénea. Al acercarse al duramen, las divisiones de la veta se vuelven bien definidas. La larga experiencia demuestra sin lugar a dudas que la madera de veta bien definida produce un mejor tono simple, e inconmensurablemente un mejor tono de doble cuerda. La experiencia también demuestra sin lugar a dudas que el duramen posee una mayor elasticidad. Por lo tanto, cuando se dobla con fuerza, al soltar la fuerza, el duramen vuelve al punto de reposo con mayor rapidez. El punto de valor en el arco indio reside en la rapidez con la que vuelve al punto de reposo. Por lo tanto, para estos arcos se selecciona la madera del duramen. De jóvenes, muchos aprendimos lo insatisfactorio que resulta un arco hecho de un árbol inmaduro o de la albura de un árbol maduro. Esta insatisfacción provenía de la lentitud con la que el arco volvía al punto de reposo. La acción de dicho arco es débil y la flecha solo puede proyectarse de forma débil. En la caja de resonancia de madera inmadura encuentro precisamente una debilidad similar; también, debilidad en el tono de las cajas de resonancia de árboles maduros que no poseen un sonido bien definido.
grano.
En mis cincuenta años de trabajo restaurando violines usados, he examinado minuciosamente muchas variedades de madera. Al realizar este trabajo, he abierto cientos de violines y ahora estoy convencido de que este método ofrece una evidencia más satisfactoria sobre las peculiaridades del sonido que cualquier otro método posible. Creo que este método revela indudablemente la razón por la que muchos artesanos hábiles no logran obtener un mayor número de obras maestras del arte del sonido. Dejando de lado el modelo y la mano de obra, ahora creo que la calidad de la madera de la caja de resonancia es fundamental para que un violín sea una obra maestra del arte del sonido. Con "obra maestra del arte del sonido" me refiero al máximo nivel de belleza sonora. No basta con que el sonido de una sola cuerda sea bello; todos los sonidos de doble cuerda también deben serlo. Es en este último aspecto donde la caja de resonancia suele demostrar su falta de valor. Si tienen la paciencia suficiente con mi lentitud, con el tiempo comprenderán mis razones para atribuir valor sonoro a la caja de resonancia. No pretendo ser infalible por tales razones. Solo alegaré que tales razones son satisfactorias para
mí mismo.
En el transcurso de mi experiencia con violines usados, he visto cambios en las superficies internas de la tapa armónica que, según creo, no son generalmente conocidos. Por ejemplo, no he oído ni leído sobre una contracción desigual de las fibras de la tapa armónica después de que el violín haya salido de las manos del luthier. Sin embargo, he visto casos, en violines de mano de obra impecable, en los que se produce una contracción desigual de la tapa armónica.
causó graves daños al tono. Recuerdo un violín con un historial impecable de 150 años, un violín indudablemente construido por un artesano, un violín en el que se encontraba una contracción desigual debajo de la cuerda Mi. Estas desigualdades de grosor se extendían a lo largo de la fibra en forma de crestas y valles, y se debían a la contracción de ciertas fibras en mayor grado que las fibras vecinas. Había tres de estas crestas y dos valles, y su ubicación, por sí sola, podía dañar el tono de la cuerda Mi. El único remedio posible consistía en restablecer un grosor uniforme en esa zona en particular. Les aseguro que no fue necesario ningún otro trabajo, de ningún tipo, ni en el interior ni en el exterior de este violín. El resultado de restablecer la uniformidad del grosor en esta zona fue la eliminación del "ruido" en el tono de la cuerda Mi. Estoy convencido de que el trabajo de graduación original en esta caja de resonancia se realizó con precisión.
[Aunque no viene al caso, lucha contra mi naturaleza omitir mencionar que la dulzura cautivadora del tono que ahora posee este viejo violín es algo que se debe conservar en la memoria. Su potencia tonal solo se ve ligeramente disminuida por la edad y el uso. El grosor de la caja de resonancia se midió con tanta precisión para la cuerda Mi de calibre 2 que la pequeña cantidad de madera que se quitó debajo de esa cuerda resultó en una disminución perceptible de su potencia tonal. Un ejemplo de contracción desigual en la caja de resonancia, después de salir de las manos del constructor, puede ser una explicación suficiente de una de las razones por las que algunos violines bien hechos, que poseen
Madera superior para la caja de resonancia, no logra convertirse en obras maestras del arte tonal. Sin embargo, ofreceré otro ejemplo de contracción en la madera de la caja de resonancia, quizás aún más sorprendente. Así pues: De una tabla que tuve en mi poder durante cuarenta años, fabriqué una caja de resonancia. Por ciertas razones, esta caja de resonancia no se utilizó hasta dos años después. Mi sorpresa fue enorme al descubrir una contracción adicional posterior a su graduación. Este hecho demuestra que la cantidad modifica considerablemente la contracción de la madera durante el secado. La viruta fina de la tabla, por muy seca que esté, puede secarse aún más. La viruta fina también se hincha más rápidamente en presencia de calor y humedad. Estos hechos son de interés para el estudiante de violín. Por lo tanto, no importa cuánto tiempo haya estado secándose el bloque o la tabla hendida, debemos considerarla, en cierto modo, como madera nueva; y, cuando la caja de resonancia, hecha de un bloque o una tabla, se reduce a grosores considerados como el límite de seguridad, no debemos sorprendernos de una mayor contracción. Tampoco debemos sorprendernos de la mayor rapidez de la hinchazón en presencia de calor y humedad. El calor y la humedad, como se demostrará más adelante,
Más tarde, son enemigos implacables del violín. También se unen para frustrar la intención del violín.
constructor.]
Debido a que el vapor de agua puede disminuir considerablemente tanto la resonancia como el brillo del tono, dependo de un higrómetro para registrar el porcentaje de saturación existente en el aire de mi taller, incluso al reafinar violines usados durante mucho tiempo. A medida que la humedad en
La madera se seca con el aire seco, por lo que, cuando es necesario, utilizo calor artificial para asegurar su sequedad. Una vez que me aseguro de que se ha eliminado toda la humedad de la madera, procedo al sellado hermético que describiré más adelante.
La contracción desigual no ocurre en todas las tapas armónicas. No sé cómo predeterminar la contracción desigual. Sin duda, el constructor de este violín jamás imaginó que la futura contracción de algunas fibras en la tapa armónica perjudicaría el sonido. Considero que la mano de obra exhibida en el interior de este violín alcanza el límite de la habilidad humana. Que su tapa armónica es una buena selección, aparte de esas pocas fibras encogidas, podemos saberlo por el hermoso sonido. El siguiente punto también contiene algo de interés para el estudiante de violín, es decir, el sonido defectuoso de este violín nunca podría haberse mejorado con el uso del arco. A pesar de la antigüedad y la mano de obra de este violín, debido a defectos de sonido por causas imprevistas, su valor era insignificante. Hoy su valor ha cambiado. Quien puede comprar este violín ahora, ya es rico. ¿Dónde reside su valor? No está en la potencia del sonido. Su sonido no es ni potente ni débil. La pregunta es difícil de responder. No está en la antigüedad. Hay violines de mayor antigüedad, pero de menor valor. No se trata de la dulzura del tono. Hay violines de tono igualmente dulce, pero de menor valor. La respuesta desafía la enunciación. De manera inadecuada, la pregunta se responde diciendo: "Su tono despierta una sensación de ten-
"Los sentimientos tiernos nos hacen desprendernos de la riqueza. No conozco otra lógica que explique por qué alguien pagaría más por un violín de dulce sonido que por otro de igual dulzura. Por lo tanto, el precio pagado puede considerarse una medida de la ternura del comprador. Pero, por qué, o cómo, un violín de dulce sonido despierta mayor ternura que otro, sigue siendo para mí un enigma sin resolver. El hecho es que existe. Eso es todo lo que sé al respecto."
Sin embargo, hay algo más que se podría decir al respecto. Así pues: Cuando el sonido de un violín despierta sentimientos tiernos en el intérprete, entonces, y por ello, el intérprete nos brindará su máxima expresión musical. Al vender un violín de este tipo,
El dueño le pone precio a sus sentimientos.
En el curso de mi trabajo con violines usados, me familiaricé con maderas de caja de resonancia de diferentes países, como Escandinavia, Rusia, Francia, Suiza, Tirol austríaco e Italia. En resumen, la madera de caja de resonancia más resonante que observé crecía en el Tirol austríaco; sin embargo, algunas muestras de Suiza e Italia han resultado difíciles de superar en esta importante cualidad. Sin duda, el calor y la humedad modifican en gran medida la uniformidad del crecimiento anual, la densidad de la fibra y la presencia o ausencia de grasa en el género Pinus. Debido a que la altitud modifica la cantidad de calor y humedad, se deduce que el pino, con anchos de grano muy variables y cantidades variables de grasa, puede encontrarse en todas las localidades montañosas donde crece el pino. Por lo tanto, cada muestra
Ningún árbol de pino, de ningún país, sirve como caja de resonancia. Para tal fin, el pino presenta cinco defectos graves; y encontrar un solo árbol sin alguno de estos defectos, incluso en las zonas más privilegiadas, es una tarea difícil. Estos cinco defectos son:
(1) Irregularidad en el ancho del grano.
(2) Grano torcido.
(3) Grano sin marcar.
(4) Demasiada densidad de grano.
(5) Grasa o tono.
De algunas zonas de Noruega procede pino de veta perfectamente recta y uniforme; pero, si bien posee una sonoridad notable, su densidad es tan elevada que le confiere un timbre algo desagradable y punzante. Si su densidad fuera menor, creo que no se podría encontrar mejor madera para la caja de resonancia.
Suponiendo que usted esté familiarizado con la madera para cajas de resonancia cultivada en el extranjero, paso a considerar la madera nacional. Los extranjeros afirman que nuestra madera no tiene valor para su uso como caja de resonancia. Se equivocan. Colorado, en cierta medida, ofrece abeto de valor insuperable para este fin; mientras que Michigan, también en cierta medida, ofrece pino muy superior a la mayoría de las cajas de resonancia europeas "prefabricadas" que amablemente nos envían. El tipo de pino de Michigan al que me refiero es peculiar por el color de su veta. Así, en el pino europeo las diferentes vetas están separadas por líneas oscuras; pero las vetas de este pino de Michigan están separadas por líneas blancas. Otra peculiaridad de este pino de Michigan es el hecho de que su
La línea blanca es la parte más suave de la veta. Hay otra peculiaridad aún más llamativa. Posee, como marcas transversales, numerosas escalas de brillo plateado; y, como sé por observación, estas marcas transversales conservan su brillo durante más de medio siglo. Excepto por la mera potencia del tono, esta calidad de pino doméstico, para su uso en la caja de resonancia, es insuperable por cualquier otro pino. El tono de esta madera, para estudio, salón o salas de conciertos más pequeñas, no puede ser superado hasta donde alcanza mi observación. Incluso en violines nuevos, los tonos de doble cuerda de esta madera poseen un atractivo extraordinario. Su valor solo carece de gran potencia. Sobre esta madera, el barniz debe aplicarse con moderación. No tiene aspereza de tono que deba ser suprimida. Para la barra y el alma, no he usado otra madera en muchos años. Como barra, de ninguna otra madera puedo obtener la misma belleza para el tono de la cuerda G.
Mi experiencia con el cedro blanco para la fabricación de cajas de resonancia es bastante limitada, ya que solo lo he probado dos veces. Para estas pruebas utilicé cedro que había secado al aire libre durante más de treinta años bajo mi propia observación. Una peculiaridad notable de esta madera es su capacidad para resistir la descomposición por calor y humedad. Incluso después de haber estado expuesta durante tantos años, su superficie solo mostraba leves rastros de pudrición, mientras que debajo, el color era sumamente brillante. La presencia de grietas profundas dificultó mucho la tarea de asegurar las muestras para la fabricación de cajas de resonancia. De hecho, solo uniendo cuatro piezas pude obtener una caja de resonancia perfecta.
Muestra de madera de grano recto y de buena calidad. Su rápido crecimiento produce una veta muy ancha. En mi opinión, por lo general, la anchura de la veta es excesiva para una caja de resonancia de alta calidad. Como mera suposición, su anchura de veta podría no ser un problema para la caja de resonancia del bajo. El timbre de esta madera, según mis dos experimentos, era peculiar. El volumen del sonido era muy marcado, pero la intensidad muy débil. Así, a corta distancia, el sonido era potente; a larga distancia, no se podía forzar su sonido.
La intensidad del sonido es algo que desafía toda explicación. Sabemos que el sonido puede poseer características distintivas.
El volumen puede variar, aunque solo sea capaz de propagarse a distancias cortas; y, por otro lado, el sonido puede tener un volumen disminuido, aunque sea capaz de propagarse a distancias mayores. Sin duda, la causa de este fenómeno reside en la peculiaridad del agente productor de sonido. La voz humana ofrece abundantes ejemplos de las diferencias entre volumen e intensidad del sonido. Algunas voces pueden propagarse fácilmente a distancias imposibles para otras; y las voces que viajan a mayor distancia pueden no sonar más de la mitad de fuertes que las voces cercanas. El mismo fenómeno se presenta en la caja de resonancia. La calidad de la intensidad no puede predeterminarse en ningún agente productor de sonido. Una marcada intensidad de tono es una valiosa característica del violín. Es una característica que influye en gran medida en los precios. Como mera propuesta comercial para el solista, el valor del violín debería basarse en la calidad del sonido.
El tono del violín A satisface al oyente a 400 pies. Empleando 100 como unidad, el cuadrado de 4 unidades es igual a 16. El tono del violín B satisface a 800 pies; y el cuadrado de 8 unidades, igual a 64, por lo tanto, el valor real y práctico de estos dos violines es de 16 a 64. El tono de estas tapas de cedro blanco era de una calidad desagradable y áspera. Curiosamente, algunos violinistas consideraron que dicho tono áspero era bastante satisfactorio. Se requirió una capa gruesa de barniz para atenuar la aspereza del tono; pero, al eliminar la aspereza, se produjo una pérdida notable de potencia tonal. Había aún una característica peculiar en el tono de estas tapas. Así pues: una fuerte presión del arco no producía un tono mayor que una presión moderada del arco. Por estas razones, no considero que el cedro blanco tenga el mismo valor tonal que el pino de Michigan aquí descrito.
Que la caja de resonancia del violín sea responsable del valor tonal del violín es una cuestión que ya no está en discusión. Por lo tanto, para el estudiante de violín, la madera para la caja de resonancia se convierte en una cuestión de suma importancia. Aparentemente, una vida dedicada a reajustar la afinación de violines usados podría permitir predeterminar con precisión la calidad tonal de muchos grados de pino y abeto. Aunque tal experiencia ofrece bases para conjeturas cercanas, les aseguro mi incredulidad en la posibilidad de predeterminar con precisión las cualidades tonales de cualquier muestra dada de madera para la caja de resonancia. En mi experiencia, las sorpresas que aguardan a la prueba del arco han sido infinitas. Describiré un grado de pino pos-
evaluando la aproximación más cercana a la uniformidad de la calidad del tono. Pero es necesario tener en cuenta que mi descripción se basa en las cualidades físicas y la apariencia del pino curado durante mucho tiempo después de haber sido moldeado en forma de caja de resonancia. Esta necesidad surge del hecho de que la rapidez de los cambios de color se ve muy modificada por la cantidad de madera. Así, el brillo original del color se conserva más tiempo en el tronco que en el bloque hendido; también se conserva más tiempo en el bloque que en la caja de resonancia. Por lo tanto, el color no puede ser una evidencia confiable para determinar el tiempo transcurrido desde que se cortó una muestra de madera determinada del tocón. Y además, diferentes muestras de pino para cajas de resonancia adquieren profundidades de color muy diferentes con el avance de la edad. Tengo una caja de resonancia, de edad respetable, que muestra apenas una sombra de cambio de color bajo la superficie. En otras, el cambio de color es marcado, habiendo pasado a un marrón rojizo por completo. La causa de los cambios de color que se producen durante el tiempo de curado es un misterio que no tiene explicación. Pero, por larga observación, sé que diferentes cualidades tonales siguen diferentes profundidades decoloreando la mesa de resonancia.[*]Así, desde la caja de resonancia que conserva durante más tiempo su brillo original, el tono tiene una cualidad algo punzante que exige una arcada sumamente cuidadosa; mientras que, desde la madera marrón rojiza, el tono es bastante agradable independientemente de la arcada. La mayor diferencia se muestra en los tonos de doble cuerda. Desde la madera clara, el efecto más bello de los tonos de doble cuerda es imposible, mientras que desde la madera marrón rojiza,
Los tonos de doble cuerda poseen el grado más alto de agradable. Que los hermosos tonos de doble cuerda, nada del violín proporciona mayor placer al oyente. Del violinista solista, los tonos de doble cuerda son imperativamente exigidos. Por lo tanto, al seleccionar un violín para uso solista, este conocimiento se vuelve valioso. Debido a que la tapa armónica marrón rojiza se agrieta fácilmente, podemos suponer que el tejido conectivo entre las fibras más densas es de débil
poder. También podemos suponer que la ausencia de
La fuerza del tejido conectivo permite una mayor independencia en la acción de las fibras, y a esta independencia se deben los agradables tonos de doble cuerda. La madera clara no se parte fácilmente. Su tejido conectivo es mucho más resistente. Por lo tanto, sus fibras, o la parte más densa de la veta, tienen menor independencia en su acción.
Para ayudar a explicar el color marrón rojizo del pino curado, ofrezco lo siguiente:
Allá por los años cincuenta del siglo pasado, un granero y un cobertizo, de varios cientos de pies de largo, estaban revestidos con pino de Michigan. En aquellos tiempos, solo se talaban grandes árboles de los pinares para obtener madera. La sierra de guillotina antigua estaba en su máximo esplendor. Esos grandes troncos se cortaban de manera transversal, y cada tabla o tablón se canteaba con una sierra más pequeña. Así, la madera era de varios anchos y, sin clasificarla como en la actualidad, se vendía al consumidor. Gran parte de la madera que cubre este granero en particular se clasificaría hoy como "de primera calidad".
Las tablas tenían 30 pulgadas de ancho. "Vestidas"
La madera era entonces desconocida. Casi también la pintura. El revestimiento de este granero no estaba ni "preparado" ni pintado jamás. Cuarenta y cinco años después examiné estas tablas sin protección. La mayoría de las tablas más anchas estaban partidas por la mitad. Pude determinar fácilmente cada tabla que provenía del centro del tronco por la evidencia de desintegración. Así: En la superficie de todas esas tablas "cortadas del centro", la fibra conectiva blanda se ha descompuesto, convertido en polvo y descompuesto a diferentes profundidades en diferentes tablas, o mejor dicho, en tablas de diferentes troncos. La parte más densa de la veta permaneció como crestas. [Algunos violines viejos y desgastados presentan crestas similares en los extremos de la caja de resonancia; y otros no. La diferencia es claramente atribuible a diferentes grados de dureza del tejido conectivo. En mi observación, esas cajas de resonancia viejas, que mostraban una mayor pérdida de tejido conectivo, poseían invariablemente un mayor valor tonal en doble cuerda y una mayor suavidad tonal en general. Además, este tipo de cajas de resonancia se parten con mayor facilidad.
En todas esas tablas desgastadas por el clima provenientes de esa parte del tronco cercana a la corteza, las crestas de fibra más densa son mucho menos prominentes. De hecho, algunas de estas tablas no muestran ninguna cresta. Una peculiaridad llamativa de estas crestas es su línea. Esta línea, en la gran mayoría de esas tablas viejas, sigue un curso en zigzag. Solo en muy pocas de ellas la línea de la cresta es recta. Ahora llego a la cuestión del color. Al preparar la superficie expuesta al clima de todas las tablas "de corte central", aparece una gran variedad de colores. Pero, solo en la más ancha
En las tablas, encuentro el color marrón rojizo. Sin duda, esas tablas de 76 centímetros (30 pulgadas) provenían de árboles maduros. Ahora vuelvo a las tablas que muestran vetas rectas. Su color varía del pálido al color mantequilla, y su efecto cromático general es brillante. En muy pocas de ellas encuentro marcas transversales plateadas brillantes, o escamas.
La evidencia proporcionada por estas tablas antiguas apunta a la edad del árbol en pie como la fuente del color marrón rojizo desarrollado por el secado. En la carpintería hay mucho valor para el estudiante de violín. En ella podemos aprender que el tejido conectivo en los árboles jóvenes o inmaduros es mucho más resistente que en el árbol maduro; también, que la madera inmadura, si bien permite mayores grados de flexión, al soltarla permanece parcialmente doblada; mientras que la madera madura, si bien permite menos flexión, al soltarla regresa rápidamente a su punto de reposo original. Esta característica de regresar rápidamente al punto de reposo es de gran importancia para los agentes productores de tono. La experiencia demuestra claramente que la madera inmadura, para su uso en la caja de resonancia, es inferior a la madera madura. En mi experiencia en el reafinado de violines, he encontrado cajas de resonancia hechas de dos muestras de madera tomadas de árboles de edades muy diferentes. Dichas cajas de resonancia han demostrado invariablemente que la madera madura posee un mayor valor tonal. Así, en tales cajas de resonancia, cuando la mitad izquierda es de madera madura, mientras que
La mitad derecha es de madera inmadura, las cuerdas G y D producirán un buen tono musical, mientras que A y E
Las cuerdas producirán un tono "muerto" o sin vida;
y viceversa, los tonos G y D carecerán de vida.
He llegado a un punto sobre el mecanismo de la caja de resonancia del violín que parece no haber recibido mucha atención.
Me refiero a la acción independiente de fibras contiguas.
Desde mi punto de vista, y siendo correctos los demás datos, la acción independiente de las fibras contiguas en la zona de resonancia del violín determina su valor tonal. Esta opinión se basa en la observación de que una madera con tejido conectivo rígido no puede producir un sonido agradable. He constatado este hecho en numerosos casos. Las siguientes consideraciones explican la necesidad de la acción independiente de las fibras en la caja de resonancia: El tono La requiere que ciertas fibras de la caja de resonancia golpeen el aire contenido 450 veces por segundo; el tono Mi requiere 675 golpes. Para la producción simultánea de los tonos La y Mi, las fibras contiguas de la caja de resonancia deben vibrar a diferentes frecuencias por segundo. La diferencia para estos tonos es la diferencia entre 450 y 675. Por lo tanto, las fibras que producen Mi vibran 225 veces más por segundo que las que producen La. Que las fibras que producen los tonos La y Mi son contiguas se demostrará más adelante. En este caso, la evidencia parece concluyente. Nuevamente llamo su atención sobre el pino que posee el color marrón rojizo como un grado que ofrece la mínima rigidez del tejido conectivo. Como evidencia adicional relacionada con la cuestión de la acción independiente de
Con fibras contiguas, les presento este violín que, según mi conocimiento, ha estado en uso durante 40 años. Dos veces lo he abierto en un intento por lograr que sus tonos de doble cuerda fueran agradables. Al aplicar un arco, se percibe de inmediato que mis esfuerzos han fracasado. Declaro que obtener tonos de doble cuerda agradables de esta caja de resonancia es imposible; y la razón de tal imposibilidad se debe a la rigidez inherente del tejido conectivo. Este violín tiene cierta historia; y, dado que la historia de algunos violines es el único valor que poseen, presento la historia de este como prueba. Al abrirlo por primera vez, encontré bajo una acumulación de suciedad una etiqueta legendaria con la siguiente inscripción:
Andrés Guamerius
Subtitulo Santa Thtresia 1645. En apariencia, esta etiqueta es la encarnación de la inocencia; y, en contorno y perfil, este violín es una réplica exacta del Andreas. En mis observaciones anteriores del violín, este se recuerda como el que daba
orgullo creciente para el Prof. . Muchas veces tuve
Lo contempló con asombro boquiabierto. ¿Quería poseerlo? ¡Claro que sí! Pero yo no tenía la riqueza de Creso. Otro joven fue más afortunado. La fortuna sonrió a este joven en particular.
compañero. El profesor encontró la necesidad de mudarse a
otras partes. Pero, antes de que pudiera hacer tal movimiento, su irrazonable casera le exigió el pago de su pensión. Bajo tales circunstancias angustiosas, el profesor, a su pesar, "bajó" su precio,
y que el joven Creso se convirtió en el orgulloso propietario.
de un Andreas Guarnerius. El tiempo no solo pasa volando, sino que también trae consigo cambios vertiginosos. Este joven, obligado a ganarse la vida, guardó su Andreas en el desván junto con su otro violín de valor común. Los ratones siempre prefieren vivir en el desván. Es la habitación más cálida de la casa. Hasta entonces, no se consideraba a los ratones como "expertos" en violines. Pero ahora, la prueba de su juicio experto sobre los violines queda demostrada por el hecho de que roeron el Guarnerius. Por eso este violín llegó a mí para ser reparado.
[Aunque no soy un experto en violines antiguos, en lo que respecta al violín moderno, en los 40 años transcurridos desde que vi este por primera vez, he aprendido algo. Por lo tanto, no me sorprende leer en la superficie interior de esta caja de resonancia, y escritas con lápiz sobre la madera, las siguientes palabras:
Johann Winkerline, Mittenwald El primer mes de octubre, 1853.
[La sugerencia de Mittenwald deja bastante claro que Andreas murió algún tiempo antes de fabricar este violín.]
Reduje el grosor de la caja de resonancia hasta lo que considero el límite de seguridad, pero solo logré eliminar ligeramente su timbre amaderado. Tras un tiempo de uso, volví a abrir este instrumento de tortura y le di a la caja de resonancia una hora de masaje.
[Como saben, hoy en día el masaje se aplica con bastante frecuencia a personas que sirven de caja de resonancia a cambio de una tarifa. Cuando dicha tarifa es elevada, hay poca diferencia en el
tono de estas dos variedades.]
Este tratamiento sí que disminuyó perceptiblemente la desagradable sonoridad de los tonos de doble parada, pero, en una habitación pequeña y vacía, siguen sugiriendo autodestrucción.
Creo que mi experiencia me permite afirmar categóricamente que, en violines con cajas de resonancia de una rigidez y densidad similares a las de esta muestra, la desesperación y la muerte preceden a la mejora del tono.
«Doctor, ¿quiere decir que los violines no siempre mejoran con la edad y el uso?»
Lo que quiero decir es que una larga vida no es suficiente para apreciar una mejora en el sonido de algunos violines.
"¿Cómo podemos entonces seleccionar violines nuevos con la certeza de una mejora en el sonido?"
Mediante la capacidad de juzgar correctamente la madera de la caja de resonancia y de evaluar su graduación a partir del tono que produce.
"Pero, doctor, ¿todo esto requiere experiencia?"
Ciertamente, señor. Pero a falta de experiencia es bastante seguro confiar en la experiencia de algún hábil,
ambicioso, concienzudo, conocedor del tono, violinista, amante del violín, fabricante de violines. Tienes mi
garantía de que tales fabricantes de violines aún no lo son
todo en el cielo.
